% Fresh phase-three supplement; phase-two artifacts are not modified.
\section{Combining equations over number fields by division norms}
\label{sec:division-norm-compiler-phase3}

This section concerns the faithful encoding of polynomial systems.
It does not assert that extending the arithmetic source to a number
field defines that field's ring of integers.  The distinction is
essential: a correct equation-combining operation preserves whatever
solutions the extended source has, including any arithmetic defects.

\subsection{A rational norm form adapted to a specified number field}

The reduced norm of a central division algebra of degree $d$ is a
homogeneous form of degree $d$ in $d^2$ coordinates, vanishing only
at the origin.  We give an explicit cyclic construction and a local
proof of division, so that the effectivity and the coefficients are
specified.  Standard background on cyclic algebras is available in
\cite[Chapter IV]{MilneCFT2020}; the prime-selection step uses
Chebotarev's theorem, for example \cite[Section 2.4]{KedlayaCFT}.

\begin{theorem}\label{thm:division-norm-conjunction-phase3}
Let $K$ be any specified number field and let $r\ge1$.
Put $d=\lceil\sqrt r\rceil$.  There is an effectively specified
homogeneous polynomial
\[
                 N\in\mathbb Z[Z_1,\ldots,Z_{d^2}]
\]
of degree $d$ such that, for $z_j\in K$,
\[
                  N(z_1,\ldots,z_{d^2})=0
                  \quad\Longleftrightarrow\quad z_1=\cdots=z_{d^2}=0.
\]
Consequently a system of $r$ equations over $K$, of degrees at
most $D$, can be combined into one equation of degree at most
$dD$, with no additional variables.  The combining form can have
rational integer coefficients even though its choice depends on $K$.
\end{theorem}
\begin{proof}
For $r=1$ use the equation itself.  Assume $d\ge2$.  Let $K'$ be
the normal closure of $K/\mathbb Q$.  Choose a rational prime
$q\equiv1\pmod d$ unramified in $K'$, and let $L$ be the unique
degree-$d$ subfield of $\mathbb Q(\zeta_q)$.  Such a prime exists
by Dirichlet's theorem.  Every nontrivial subfield of $L$ is
ramified at $q$, whereas $K'$ is unramified there.  Thus
$L\cap K'=\mathbb Q$.  Both extensions are Galois, so their
compositum has the product Galois group.

Chebotarev's theorem supplies a rational prime $p\ne q$ which
splits completely in $K'$ and whose Frobenius generates
$\operatorname{Gal}(L/\mathbb Q)$.  Hence $p$ is inert in $L$.
Fix a generator $\sigma$ of this cyclic Galois group and form
\[
 A=(L/\mathbb Q,\sigma,p)
   =\bigoplus_{j=0}^{d-1}L u^j,
 \qquad u^d=p,\qquad ua=\sigma(a)u.
\]
These relations give the usual associative cyclic algebra.

At a place of $K$ above $p$, the completion is $\mathbb Q_p$,
and the corresponding completion of $L$ is the unramified
degree-$d$ extension $L_p/\mathbb Q_p$.  On the local cyclic
algebra define
\[
 w\left(\sum_{j=0}^{d-1}a_j u^j\right)
       =\min_{a_j\ne0}\left(v_{L_p}(a_j)+\frac jd\right).
\]
The minimum is attained at exactly one index, since the
fractional parts $j/d$ are distinct.  Multiplying the unique
lowest terms shows that $w(ab)=w(a)+w(b)$: the automorphism
preserves the valuation, and reducing $u^d$ to $p$ preserves
the weight.  No other product term has this lowest weight.
The local algebra therefore has no zero divisors and is a
division algebra, because it is finite dimensional over a field.
It follows that $A\otimes_{\mathbb Q}K$ is also a division
algebra: a zero divisor would remain one after extension to
that completion.  The same argument proves that $A$ is division.

Here is a concrete basis of $L$ consisting of algebraic integers
for writing its norm form.
Let $g$ generate $(\mathbb Z/q\mathbb Z)^*$, let
$H=\langle g^d\rangle$, and put
\[
 \eta_i=\sum_{h\in H}\zeta_q^{g^ih},\qquad 0\le i<d.
\]
The periods are linearly independent over $\mathbb Q$.
Indeed a relation among them gives a polynomial in $\zeta_q$
of degree at most $q-1$ and constant coefficient zero.
It must be a multiple of $1+T+\cdots+T^{q-1}$, whose constant
coefficient is one, and therefore must be zero.  Since the
periods are fixed by $H$, they form a basis of $L$.  Their sum
is $-1$, so $1,\eta_1,\ldots,\eta_{d-1}$ is also a basis.
Choose
$\sigma$ so that $\sigma(\eta_i)=\eta_{i+1}$, with indices
modulo $d$.

Write $a_j=Z_{j,0}+\sum_{i=1}^{d-1}\eta_i Z_{j,i}$.  Let
$D(a)=\operatorname{diag}(a,\sigma(a),\ldots,\sigma^{d-1}(a))$
and define the matrix $U$ by
\[
 Ue_j=e_{j-1}\ (1\le j<d),\qquad Ue_0=p e_{d-1}.
\]
Then $U^d=pI$ and $UD(a)=D(\sigma(a))U$.  The norm form is
\begin{equation}\label{eq:explicit-cyclic-norm-phase3}
 N(Z)=\det\left(\sum_{j=0}^{d-1}D(a_j)U^j\right).
\end{equation}
Applying $\sigma$ to this matrix conjugates it by $U$, so
its determinant has rational coefficients.  The periods are
algebraic integers and $p$ is an integer, so those rational
coefficients are integers.  The form is homogeneous of degree
$d$.  A nonzero coordinate vector over $K$ represents a
nonzero, hence invertible, element of $A\otimes K$; its matrix
is invertible and its norm is nonzero.  This proves anisotropy.

The construction is effective.  One can search for $q$, compute
its cyclotomic periods, and search primes $p$ testing complete
splitting in $K'$ and the indicated Frobenius class in $L$.
The cited existence theorems ensure termination.  Finally
substitute the $r$ residual polynomials for the first $r$
coordinates of $N$, and zeros for all remaining coordinates.
Anisotropy proves equivalence, and homogeneity gives the degree
bound without introducing any quantified coordinate.
\end{proof}

\begin{remark}[Coefficient bound]
In every complex embedding, $|\eta_i|\le(q-1)/d$.
Every entry of the matrix in
\eqref{eq:explicit-cyclic-norm-phase3} is a linear form whose
coefficient absolute sum is at most $(q-1)p$.  Expanding the
determinant and applying the triangle inequality therefore gives
\[
                       \|N\|_1\le d!\bigl((q-1)p\bigr)^d.
\]
This includes all determinant permutations and all monomials.
The form is primitive: assigning $Z_{0,0}=1$ and all other
coordinates zero represents the identity and gives $N=1$.
If all input residuals have
coefficient norm at most $M$, their combined polynomial has
coefficient norm at most $\|N\|_1M^d$.
\end{remark}

\subsection{The preferred degree-36 form and the two new quadratic systems}

The preferred phase-three systems have respectively 1315 and 1339
total coordinates, with 1272 and 1296 quadratic residuals.
Both residual counts are at most $36^2=1296$.  For these systems
take the completely explicit choices
\[
                    d=36,\qquad q=37,\qquad p=2.
\]
The element 2 generates $(\mathbb Z/37\mathbb Z)^*$: indeed
$2^{18}\equiv-1$ and $2^{12}\equiv26\not\equiv1\pmod{37}$,
which exclude the maximal proper divisors of 36.
Thus $L=\mathbb Q(\zeta_{37})$ has degree 36 and two is inert.
The preceding local valuation proof shows that
\[
                    A=(L/\mathbb Q,\sigma,2),
             \qquad \sigma(\zeta_{37})=\zeta_{37}^2,
\]
is division.  Formula~\eqref{eq:explicit-cyclic-norm-phase3}
therefore gives an integer norm form of degree 36 in 1296
coordinates.  Here the period basis is simply
$1,\zeta_{37}^{2},\zeta_{37}^{2^2},\ldots,
\zeta_{37}^{2^{35}}$, with exponents reduced modulo 37.

\begin{lemma}\label{lem:division-coprime-base-change-phase3}
Let $A$ be a central division algebra of degree $d$ over a
field $k$.  If $K/k$ is finite and
$\gcd(d,[K:k])=1$, then $A\otimes_k K$ is division.
\end{lemma}
\begin{proof}
Write $A\otimes K\simeq M_s(B)$, where $B$ is central
division over $K$ of degree $d/s$.  A simple left module
has $K$-dimension $s(d/s)^2=d^2/s$.  By restricting scalars
it is a vector space over the division ring $A$, so its
$k$-dimension is a multiple of $d^2$.  Hence
$s\mid[K:k]$.  Since also $s\mid d$, coprimality gives
$s=1$.
\end{proof}

This proves that the fixed degree-36 form remains anisotropic
whenever $\gcd(36,[K:\mathbb Q])=1$.  The assertion must not
be justified by choosing one completion of coprime local degree:
that inference is false for composite $d$, whereas the module
argument applies directly.

The same form also remains anisotropic over $\mathbb Q(i)$.
Here is a local proof, independent of the coprime-degree lemma.
At 37 put $\pi=\zeta_{37}-1$.  The residue field of
$L_{37}$ is $\mathbb F_{37}$ and
$\sigma(\pi)/\pi\equiv2\pmod\pi$.
The algebra $\mathcal O_{L_{37}}[u]$, $u^{36}=2$, has residue
algebra
\[
                    \mathbb F_{37}[u]/(u^{36}-2),
\]
which is a field.  Indeed a root $\alpha$ satisfies
$\alpha^{37}=2\alpha$, so its Frobenius orbit has length
exactly 36 because 2 has multiplicative order 36.
For $x=\sum a_j u^j$ define $w(x)=\min_j v_\pi(a_j)$.
After factoring out its lowest power of $\pi$, a nonzero
element has nonzero reduction in this residue field.
Commuting $\pi^s$ through powers of $u$ changes reduction
by the automorphism $u\mapsto2^su$; consequently the lowest
reductions of a product are still a product of nonzero
elements.  Thus $w(xy)=w(x)+w(y)$ and the local algebra
is division.  Finally 37 splits in $\mathbb Q(i)$ because
$6^2\equiv-1\pmod{37}$, so a completion after this base
change is still $\mathbb Q_{37}$.  Division at that completion
proves division over $\mathbb Q(i)$.

Combining the two preferred quadratic systems gives polynomials
with respectively 1315 and 1339 total coordinates and
\emph{degree exactly 72}.  The exact residual inventories put
the input monomial $X_0^2$ with coefficient one in the first
residual; no other residual contains $X_0$.  The first norm
coordinate is the coefficient of the algebra identity, so the
coefficient of $X_0^{72}$ in the resulting polynomial is one.
In particular the top degree cannot cancel, and each resulting
integer polynomial is primitive.

The coefficient estimate becomes
\[
                      \|N\|_1\le36!\,72^{36}<2^{361}.
\]
The exact residual tables have coefficient norm at most
$2^{511}$.  The 1339-coordinate version, with ordinary
integer coefficients already materialized in those residuals,
therefore has coefficient height at most
\[
                              2^{18757}.
\]
For the 1315-coordinate version substitute the fixed integer
$\kappa=3^{2^{28}}$ into residuals affine in $\kappa$.
Using $\log_2\kappa<429496730$ gives coefficient height at
most
\[
                              2^{15461901037}.
\]
The artifact names
\texttt{norm\_repair\_1315\_degree72\_20261106\_v1.json.gz}
and
\texttt{norm\_repair\_1339\_degree72\_20261106\_v1.json.gz}
specify the complete residual lists, cyclotomic basis,
matrix entries, and determinant.  They are exact finite
polynomial specifications, not expanded coefficient tables
of degree 72.  The symbol $\zeta_{37}$ is eliminated by
reduction modulo its cyclotomic polynomial and extraction
of the constant coefficient; it is not a quantified variable.

For an arbitrary specified number field the fixed form may
cease to be anisotropic.  Theorem~\ref{thm:division-norm-conjunction-phase3}
provides a field-dependent choice of primes $q,p$ with the
same degree 36 and rational integer coefficients.  Thus
degree 72 and the coordinate counts persist, while the
coefficient bound becomes the explicit prime-dependent bound
in that theorem.

\subsection{The preceding degree-37 form and its comparison bounds}

The choices
\[
                         d=37,\qquad q=149,\qquad p=5
\]
are explicit.  The degree-37 subfield of $\mathbb Q(\zeta_{149})$
is unramified at five, and five is inert there: the Galois group
has prime order 37 and $5^4\equiv29\not\equiv1\pmod{149}$.
The local valuation argument above proves division.

This same rational norm form remains anisotropic over every
number field $K$ for which $37\nmid[K:\mathbb Q]$.  To see
this directly, use
$\sum_{v\mid5}e_vf_v=[K:\mathbb Q]$ to choose a place with
$37\nmid e_vf_v$.  The unramified degree-37 extension remains
a field after this local base change.  Normalizing the valuation
of $K_v$ gives $v(5)=e_v$, and the weights
\[
                       v(a_j)+\frac{e_vj}{37}
\]
again have distinct fractional parts.  The same lowest-term
argument proves that the local cyclic algebra is division.
For fields with degree divisible by 37, the field-dependent
prime choices in the theorem still supply degree 37 and rational
integer coefficients.

For comparison, there are 1322 primitive quadratic residuals in the
1367-coordinate system and 1346 in the 1391-coordinate system.
Both numbers are at most $37^2=1369$.  Thus the norm construction
combines either system into one polynomial of degree exactly 74,
with the same respective total coordinate counts.  Exactness of
the degree follows without expansion: the highest homogeneous
quadratic parts of the residuals are not all zero.  At some
rational specialization this tuple is nonzero, and anisotropy
shows that its norm is nonzero.  It is exactly the homogeneous
degree-74 part of the combined polynomial.

For the explicit $(37,149,5)$ choice,
\[
               \|N\|_1\le37!\,740^{37}<2^{496}.
\]
The exact residual tables give maximum coefficient norm at most
$2^{511}$.  Hence the ordinary integer version with 1391 total
coordinates has coefficient height at most
\[
                              2^{19403}.
\]
For the 1367-coordinate version the residuals are affine in
the fixed integer $\kappa=3^{2^{28}}$.  Since
$\log_2\kappa<(8/5)2^{28}<429496730$, a valid integer
coefficient-height bound is
\[
                              2^{15891398413}.
\]
These bounds concern the norm-compiled polynomials.  Their
degree-74 determinant specifications are separate from the
frozen degree-four sum-of-squares tables.

\subsection{The finite-algebra compiler over every number field}

All algebraic steps in the guarded finite-algebra theorem work
over an arbitrary characteristic-zero field: the finite
separating-parameter bound, simple rational eigenline,
trace-adjugate identities, and reciprocal guard are unchanged.
Only the final conjunction by a sum of squares requires an
ordering.  Over a number field with a real embedding it remains
valid as written.  Over an arbitrary number field replace each
branch's sum of squares by the norm form of
Theorem~\ref{thm:division-norm-conjunction-phase3}.

If $n$ bounds the algebra rank and $c$ bounds the number of
closed residuals, a branch contains at most $n+c$ residuals:
one reciprocal residual, at most $n-1$ high-coefficient
conditions, and at most $c$ reversed traces.  Use one norm form
of degree
\[
                         d=\lceil\sqrt{n+c}\rceil
\]
and pad shorter lists with zero.  This introduces no new
coordinate.  Relative to the recorded sum-of-squares degree
bound, multiply each branch bound by $d/2$ before multiplying
over its finite indices.  This operation preserves the rational
projection over the specified number field of the guarded
finite source itself.

The exact recorded integer bounds give the following convenient
rounded estimates.  All counts include the input.
\[
\begin{array}{c|c|r|c}
\text{source}&\text{total coordinates}&d&\text{degree bound over }K\\\hline
\text{phase-three finite anchor}&41&16385&2^{9566686}\\
\text{phase-three semilocal}&43&23171&2^{250}\\
\text{phase-two finite anchor}&43&23171&2^{9566693}\\
\text{phase-two semilocal}&45&32769&2^{257}
\end{array}
\]
The combining form has integer coefficients chosen using $K$;
the number of coordinates and the displayed degree bound do not
depend on $[K:\mathbb Q]$.  These are field-adapted encodings
of the corresponding source systems.  They do not establish
that those systems define $\mathcal O_K$.
